% Source: https://artofproblemsolving.com/community/c3h2304304_mathcounts_problem
% Stable identity: aops-topic-2304304-q01
Use coordinates with $A=(0,0,0)$ and $B=(2,2,2)$, where a move adds $1$ to one coordinate. The lattice points of the large cube are the points with coordinates in $\{0,1,2\}$; all of them lie on the surface except the centre $(1,1,1)$.

\begin{asy}
usepackage("amsmath");
size(5cm);
defaultpen(fontsize(10pt));
pair P(real x,real y,real z){return (x+0.55*y,z+0.38*y);}
pen g=gray(0.55);
for(int i=0;i<=2;++i){draw(P(i,0,0)--P(i,0,2),g); draw(P(0,0,i)--P(2,0,i),g);}
for(int i=0;i<=2;++i){draw(P(0,i,2)--P(2,i,2),g); draw(P(i,0,2)--P(i,2,2),g);}
for(int i=0;i<=2;++i){draw(P(2,i,0)--P(2,i,2),g); draw(P(2,0,i)--P(2,2,i),g);}
draw(P(0,0,0)--P(2,0,0)--P(2,0,2)--P(0,0,2)--cycle,linewidth(0.9));
draw(P(0,0,2)--P(0,2,2)--P(2,2,2)--P(2,0,2),linewidth(0.9));
draw(P(2,0,0)--P(2,2,0)--P(2,2,2),linewidth(0.9));
draw(P(0,0,0)--P(0,2,0)--P(2,2,0),dashed+g); draw(P(0,2,0)--P(0,2,2),dashed+g);
draw(P(0,0,0)--P(1,0,0)--P(1,0,1)--P(2,0,1)--P(2,1,1)--P(2,1,2)--P(2,2,2),red+linewidth(1.3));
dot(P(0,0,0)); dot(P(2,2,2));
label("$A$",P(0,0,0),SW); label("$B$",P(2,2,2),NE);
draw(circle(P(1,1,1),0.06),blue);
label("$(1,1,1)$",P(1,1,1),NW,blue+fontsize(8pt));
label("right",P(1.5,0,0),S,fontsize(8pt)); label("up",P(0,0,1.5),W,fontsize(8pt)); label("back",P(2,1.2,0),SE,fontsize(8pt));
\end{asy}

\textbf{(b) first.} The number $\dbinom62\dbinom42=\dfrac{6!}{2!\,2!\,2!}=90$ counts all orders of the six moves, i.e. all monotone lattice paths from $A$ to $B$ through the $3\times3\times3$ grid of points, including those that go through the centre $(1,1,1)$, inside the cube. Those are not allowed.

\textbf{(a)} A path through the centre consists of a path from $A$ to $(1,1,1)$ (one move in each direction, $3!=6$ orders) followed by a path from $(1,1,1)$ to $B$ (again $6$ orders): $36$ paths. Every other monotone path avoids the only interior lattice point, and each of its moves joins two surface points along an edge on a face, so it runs on the surface. Hence the number of short paths is
\[90-36=\boxed{54}.\]
(This agrees with counting by the point reached after three moves: it must have coordinates $2,1,0$ in some order, $6$ choices, with $3$ ways to reach it and $3$ ways to continue, $6\cdot3\cdot3=54$.)
